\documentclass[11pt,a4paper]{article}
\usepackage[margin=26mm]{geometry}
\usepackage[T1]{fontenc}
\usepackage{lmodern,microtype,amsmath,amssymb,amsthm,mathtools}
\usepackage{xcolor}
\usepackage[colorlinks=true,linkcolor=blue!45!black,citecolor=blue!45!black,urlcolor=blue!45!black]{hyperref}
\setlength{\parskip}{2pt}
\newtheorem{theorem}{Theorem}[section]
\newtheorem{lemma}[theorem]{Lemma}
\newtheorem{corollary}[theorem]{Corollary}
\theoremstyle{remark}
\newtheorem{remark}[theorem]{Remark}
\newcommand{\Z}{\mathbf Z}
\newcommand{\Q}{\mathbf Q}
\newcommand{\R}{\mathbf R}
\newcommand{\C}{\mathbf C}
\newcommand{\F}{\mathbf F}
\newcommand{\OO}{\mathcal O}
\newcommand{\Spec}{\operatorname{Spec}}
\newcommand{\Pic}{\operatorname{Pic}}
\newcommand{\Cl}{\operatorname{Cl}}
\newcommand{\Frac}{\operatorname{Frac}}
\newcommand{\divisor}{\operatorname{div}}
\newcommand{\HD}{H_{\mathcal D,\mathrm{an}}}
\begin{document}
\begin{center}
{\Large Rational Gersten injectivity in dimension two}\\[5pt]
{\normalsize A quartic, a tame family, and an explicit unit class}\\[6pt]
{\small Independent research draft --- 7 October 2026}
\end{center}
\noindent\fcolorbox{black!35}{black!3}{%
\parbox{\dimexpr\textwidth-2\fboxsep-2\fboxrule\relax}{\small
\textbf{Review status.} This is a new proposed argument under internal
review. The cited theorems establish its inputs, not the new conclusion.
This note concerns uncompleted arithmetic local rings. It uses no
claim about a completed ring, continuous cyclic homology, or crystalline
comparison.}}

\paragraph{Scope of the constructions.}
Sections~1--6 prove the quartic statement, and
Section~\ref{sec:twodimexplicitgeometric} gives an explicit unit triple.
Section~\ref{sec:twodimtamefamily} treats every residue prime and
unbounded ranks while keeping the residue field $\F_p(T)$.
Section~\ref{sec:twodimstrictcoefficientinfinite} gives one fixed
uncompleted ring at each prime with infinite rational generic kernels
in every odd degree at least three.

\paragraph{Related recent preprints.}
The preprints \cite{PriorIntegralK3,PriorIntegralK5} claim integral
counterexamples in degrees three and five on different ramified
rings with finite residue fields. Both explicitly leave the order
of their constructed class, and its image after rationalization,
undetermined. The present argument requires infinite order and
uses a transcendental residue parameter in dimension two. Those
preprints are not inputs to the proof; their cited claims are not
treated here as externally verified. No priority claim is made.

\section{The source and an actual supported class}

Write
\[
 P(y)=\Phi_5(1+y)=5+10y+10y^2+5y^3+y^4,
 \qquad F=P(y)+x^4+Txy^3.
\]
Let $R_0=\Z_{(5)}[T,x,y]/(F)$, let
$\mathfrak m=(5,x,y)$, and put $A=(R_0)_{\mathfrak m}$.
Here $T$ remains transcendental in the residue field.

\begin{theorem}[Proposed dimension-two counterexample]
\label{thm:twodimgeometric}
The ring
\[
 A=\left(\Z_{(5)}[T,x,y]/
       \bigl(\Phi_5(1+y)+x^4+Txy^3\bigr)\right)_{(5,x,y)}
\]
is a regular local ring of dimension two with residue field $\F_5(T)$.
There is an actual infinite-order class $c\in K_3(A)$ whose image in
$K_3(\Frac A)$ is zero.
\end{theorem}

All $K$-groups in this note are Quillen groups. We use localization,
devissage, flat base change, and the projection formula
\cite{Quillen}. Closed-immersion pushforwards can first be taken in
$G$-theory and then identified with $K$-theory on regular schemes.

The ambient local ring
$\Z_{(5)}[T,x,y]_{(5,x,y)}$ is regular of dimension three with parameters
$5,x,y$. The equation $F$ has nonzero linear term $5$, so $A$ is
regular of dimension two. More explicitly,
\[
 5(1+2y+2y^2+y^3)=-(y^4+x^4+Txy^3),
\]
and the factor in parentheses is a unit; thus its maximal ideal is
$(x,y)$ and $5\in(x,y)^4$.

Fix a primitive fifth root $\rho$ and set
\[
 E=\Q(\rho),\qquad W=\Z_{(5)}[\rho],\qquad u=\rho-1.
\]
The ring $W$ is a DVR with residue field $\F_5$ and uniformizer $u$.
Write $\sigma_g(\rho)=\rho^g$ for $g\in\F_5^\times$.
The regular divisor $D=\Spec(A/x)$ is
\[
 A/x\simeq W[T]_{(u)},\qquad y\longleftrightarrow u.
\]

Borel's theorem and its regulator comparison give a nonzero
degree-three Borel class of $E$, with complex conjugation
$\sigma_{-1}$ acting by $-1$ on its rational regulator
\cite[Proposition~12.2]{Borel};
\cite[Section~6, especially (6.4)]{BunkeTamme}.
Quillen localization gives a surjection $K_3(W)\to K_3(E)$, since
$K_2(\F_5)=0$ \cite{QuillenFinite}. Choose an actual lift $\gamma$ of
a class with nonzero Borel regulator, and replace it by
\[
 \beta=\gamma-\sigma_{-1}\gamma\in K_3(W).
\]
Then $\sigma_{-1}\beta=-\beta$ exactly, and its Borel regulator is
nonzero. Pull $\beta$ to $D$ and define
\begin{equation}\label{eq:twodimsupported}
 c=i_*\beta\in K_3(A),\qquad i:D\hookrightarrow\Spec A.
\end{equation}
Its restriction to $A[1/x]$, and hence to $\Frac A$, is zero.
We prove that $c$ has infinite order.

\section{The ramified base change and its blow-up}

Put $B=A\otimes_{\Z_{(5)}}W$. This finite flat $A$-algebra is local:
its closed fiber is $\F_5(T)[u]/(u^4)$. Its maximal ideal is
$\mathfrak n=(u,x,y)$, and
\[
 B\simeq\bigl(V[x,y]/(F)\bigr)_{(u,x,y)},
 \qquad V=W[T]_{(u)}.
\]
Thus $V$ is a coefficient DVR with residue field $k_0=\F_5(T)$.
No perfectness of $k_0$ is assumed. The generic function field of $A$
is $\Q(x,y)$, by solving for $T$; hence $B$ is a domain, with function
field $E(x,y)$.

Set $\epsilon=5/u^4\in W^\times$. The equation
$\Phi_5(1+u)=0$ gives $\epsilon\equiv-1\pmod u$.
The projectivized tangent cone at $\mathfrak n$ is therefore the quartic
\begin{equation}\label{eq:twodimtangent}
C:\quad X^4+Y^4+TXY^3=U^4
       \quad\hbox{in }\mathbf P^2_{k_0}
       \text{ with coordinates }(U:X:Y).
\end{equation}
It is geometrically smooth. A singular point would have $U=0$;
on the chart $Y\ne0$ it would give a repeated root of
$s^4+Ts+1$. Its discriminant is $256-27T^4$, nonzero in $k_0$.
There are no possible singular points with $U=Y=0$.

\begin{lemma}[An arithmetic ordinary multiple point]
\label{lem:twodimblowup}
The ring $B$ is normal. The blow-up
$\pi:\mathcal Y=\operatorname{Bl}_{\mathfrak n}(\Spec B)\to\Spec B$
is regular, is an isomorphism away from the closed point, and has
the single reduced exceptional curve $C$. Moreover,
\[
 \OO_{\mathcal Y}(C)|_C=\OO_C(-1).
\]
For $a_g=\rho^g-1$, the strict transform of the prime
$\mathfrak p_g=(x,y-a_g)$ meets $C$ transversely at
$P_g=[1:0:g]$.
\end{lemma}
\begin{proof}
First, $B$ is regular away from its closed point. Over $E$, the
surface $F=0$ is smooth: $\partial F/\partial T=xy^3$; if $x=0$
then $P'(y)\ne0$, and if $y=0$ then $x^4=-5$ and $4x^3\ne0$.
Over $u=0$, the binary form $x^4+y^4+Txy^3$ is square-free and is
smooth away from the origin. Its gradient therefore proves regularity
at each nonclosed point over $u$. Since $B$ is a hypersurface, it is
Cohen--Macaulay; the established codimension-one regularity gives
normality.

For clarity, the three strict-transform equations in the ambient
blow-up of the regular local ring $V[x,y]_{(u,x,y)}$ are as follows:
\begin{align*}
 u\text{-chart:}\quad &x=uX,\ y=uY,\\[-3pt]
 &X^4+Y^4+TXY^3+
       \epsilon(1+2uY+2u^2Y^2+u^3Y^3)=0;\\
 x\text{-chart:}\quad &u=xU,\ y=xY,\\[-3pt]
 &1+Y^4+TY^3+
       \epsilon U^4(1+2xY+2x^2Y^2+x^3Y^3)=0;\\
 y\text{-chart:}\quad &u=yU,\ x=yX,\\[-3pt]
 &X^4+1+TX+
       \epsilon U^4(1+2y+2y^2+y^3)=0.
\end{align*}
In the second and third charts retain the ambient relations
$u=xU$ and $u=yU$. Their ambient rings are regular; for example
$u-xU$ has the parameter $u$ as its linear term at a point where
$x=U=0$. Reduction along the exceptional divisor gives the corresponding
charts of the smooth curve \eqref{eq:twodimtangent}. Hence each strict
transform has a nonzero linear term in a parameter tangent to the
ambient exceptional plane. This proves regularity along $C$, including
over the imperfect field $k_0$. These equations are the strict
transforms: division by the fourth power of the exceptional parameter
introduces no component contained in that parameter's zero divisor.

The associated graded ring is
$k_0[U,X,Y]/(X^4+Y^4+TXY^3-U^4)$. Thus the exceptional scheme is
exactly the reduced curve $C$. The blow-up tautological identity
$\mathfrak n\OO_{\mathcal Y}=\OO_{\mathcal Y}(-C)$ gives the asserted
normal line \cite[Section~31.33]{StacksBlowup}.

All root strict transforms lie in the $u$-chart. Put $b_g=a_g/u$;
these are units with reductions $g$. At $X=0$ that chart's equation is
$\prod_g(Y-b_g)=0$. Its $Y$-derivative is a unit at each root.
Near $P_g$, the strict transform of $\mathfrak p_g$ is generated by
$X$, and $u,X$ are regular parameters on $\mathcal Y$. This proves
the stated transverse intersection of multiplicity one.
\end{proof}

\section{A Picard lattice on the exceptional open}

Write $C^*=C\cap\{UY\ne0\}$, so on $U=1$ its equation is
$X^4+Y^4+TXY^3=1$, with $Y$ invertible.

\begin{lemma}\label{lem:twodimexceptionalpic}
If $\sum_g n_gP_g$ is principal as a divisor class on $C^*$, then
all four integers $n_g$ are equal.
\end{lemma}
\begin{proof}
It suffices to check after extending $\F_5$ to
$k=\overline{\F}_5$. Consider the smooth surface
\[
 M=\Spec k[X,V,Y,Y^{-1}]/\bigl(XV-(1-Y^4)\bigr),
 \qquad T=(V-X^3)/Y^3.
\]
On inverting $X$ its ring is $k[X^{\pm1},Y^{\pm1}]$, a UFD.
The four primes $(X,Y-g)$ generate $\Cl(M)=\Pic(M)$. Relations
are divisors of units $cX^rY^s$ on this open, so the sole relation
is the sum of the four primes. Hence
\[
 \Pic(M)=\Z^4/\Z(1,1,1,1).
\]
Every finite fiber $T=t$ is integral, even when singular. Under
$s=X/Y$, $w=1/Y$ it is $w^4=s^4+ts+1$, $w\ne0$.
The right side is not a square in $k(s)$: a rational square would be
a polynomial square, and comparing it with $(s^2+bs+c)^2$ forces
$b=c=0$, contrary to the constant term. Since $k$ contains $\mu_4$,
the fourth-root polynomial is irreducible. Thus each removed vertical
prime is the single principal divisor $(T-t)$, of multiplicity one.
The divisor-class localization sequence gives
\[
 \Pic(M)\xrightarrow{\sim}\Pic(M_{k(T)}).
\]
This generic fiber is $C^*_{k(T)}$, with the four displayed root
points. The claimed consequence follows.
\end{proof}

\section{Arithmetic persistence of the generic Picard group}

Let
\[
 \mathcal Q=\Spec E[T,x,y,y^{-1}]/(F).
\]
The substitution $v=x^3+Ty^3$ is an isomorphism
\begin{equation}\label{eq:twodimdanielewski}
 \mathcal Q\simeq\Spec E[x,v,y,y^{-1}]/(xv+P(y)),
 \qquad T=(v-x^3)/y^3.
\end{equation}
It is a smooth surface. Let $D_g$ denote $(x,y-a_g)$, let
$J_g=\OO_{\mathcal Q}(-D_g)$, and put
$\eta_g=1-[J_g]=[\OO_{D_g}]\in K_0(\mathcal Q)$.
Define the permitted arithmetic open
\[
 \mathcal T=\Spec B[1/u,1/y].
\]
If $S=R_0\setminus\mathfrak m$, then
$\mathcal T=\Spec S^{-1}\OO(\mathcal Q)$.

\begin{lemma}\label{lem:twodimpicinjection}
The root divisors give
\[
 \Pic(\mathcal Q)=\Z^4/\Z(1,1,1,1),
\]
and restriction $\Pic(\mathcal Q)\to\Pic(\mathcal T)$ is injective.
This Picard-group calculation is unchanged by any field extension of $E$.
\end{lemma}
\begin{proof}
Inverting $x$ in \eqref{eq:twodimdanielewski} gives the Laurent UFD
$E[x^{\pm1},y^{\pm1}]$. The primes above $x$ are the four $D_g$,
with multiplicity one. Its units are $cx^ry^s$ and $y$ was already
invertible. The standard divisor localization argument therefore
gives exactly the displayed lattice. The same calculation works over
every coefficient extension because all four distinct roots belong
to $E$.

Suppose $\sum_g n_gD_g$ becomes principal on $\mathcal T$, with
rational function $f$. Regard $f$ on the regular blow-up
$\mathcal Y$ of Lemma~\ref{lem:twodimblowup}. Its divisor there has
the form
\[
 \divisor_{\mathcal Y}(f)=\sum_g n_g\widetilde D_g
                    +\sum_j m_jV_j+mC,
\]
where every $V_j$ is a strict transform of a prime in $u=0$ or $y=0$.
Indeed, $\mathcal Y[1/u,1/y]=\mathcal T$; the only exceptional prime
is $C$. Restrict the associated line-bundle equality to $C^*$.
The $V_j$ miss this open: strict transforms of $u=0$ meet $C$ only
in $U=0$, and those of $y=0$ only in $Y=0$. The exceptional line
$\OO_C(-1)$ is trivial on $U\ne0$. Finally,
$\widetilde D_g$ restricts to $P_g$ with multiplicity one.
It follows that $\sum_g n_gP_g$ is zero in $\Pic(C^*)$.
Lemma~\ref{lem:twodimexceptionalpic} makes all $n_g$ equal, which
was already the zero relation on $\mathcal Q$.
\end{proof}

\section{Betti classes survive every permitted finite localization}

Fix a complex embedding of $E$, and let
$e_g=c_1^{\mathrm{top}}(\OO(D_g))\in H^2(\mathcal Q(\C),\R)$.
The Gysin sequence for the four disjoint affine lines $D_g$ has segment
\[
 H^1((\C^\times)^2,\R)\longrightarrow\R^4
       \longrightarrow H^2(\mathcal Q(\C),\R).
\]
The normalized angular class of $x$ has residue the diagonal vector,
while that of $y$ has residue zero. Consequently the classes $e_g$ have exactly the
diagonal relation and span an injected copy of $\R^4/\R(1,1,1,1)$.
We do not identify this subspace with all of $H^2$: there is also a
quotient arising from $H^2((\C^\times)^2,\R)$.

\begin{lemma}\label{lem:twodimbettipersistence}
For every $s\in S$, the map
\[
 H^2(\mathcal Q(\C),\R)\longrightarrow
 H^2(\mathcal Q_s(\C),\R)
\]
is injective.
\end{lemma}
\begin{proof}
Every irreducible $E$-divisor $Z$ contained in $s=0$ disappears on
$\mathcal T$. Its class in $\Pic(\mathcal Q)$ is zero by
Lemma~\ref{lem:twodimpicinjection}. Over $\overline E$, its geometric
components $Z_1,\ldots,Z_d$ form one Galois orbit. The absolute Galois action over $E$ on
$\Pic(\mathcal Q_{\overline E})=\Z^4/\Z(1,1,1,1)$ is trivial,
because the root generators are $E$-defined. Thus all $[Z_j]$ are
equal. Their sum is zero; torsion-freeness makes each $[Z_j]$ zero.
There are no inseparable multiplicities in characteristic zero.
The same is true of their complex base changes.

The kernel of degree-two restriction to the complement of a divisor
on a smooth complex surface is generated by the fundamental classes
of its irreducible components. This follows from the Gysin sequence
on its smooth disjoint locus; singular points and intersections have
real codimension four and do not change the degree-two support group.
All those classes here are the first Chern classes just shown to
vanish. This proves injectivity.
\end{proof}

\section{The multiplicative regulator and the conclusion}

Bunke--Tamme's regulator is natural and multiplicative, with analytic
Deligne target \cite[Remark~2.6 and Theorem~2.31]{BunkeTamme}.
At a point its weight-two target is $\C/\R(2)$, where
$\R(q)=(2\pi i)^q\R$. For a smooth complex surface $U$, the
holomorphic-degree filtration has $F^3=0$ literally, and the cone
description gives
\begin{equation}\label{eq:twodimdeligne}
 \HD^3(U,\R(3))=
 H^2(U,\C)/H^2(U,\R(3)).
\end{equation}
No properness or logarithmic compactification is used in this formula.
Indeed, the possible next kernel is zero because real cohomology
injects into complex cohomology.

If a point regulator is represented by $b\in\C/\R(2)$ and a line
has topological Chern class $e$, their product in
\eqref{eq:twodimdeligne} is represented by $2\pi i\,b e$.
For completeness, in the interval model of
\cite[Definition~2.7, Proposition~2.8 and Corollary~2.14]{BunkeTamme},
the point class is $-b\,dt$, and a unitary Chern form is
$2\pi i\,e$, constant in $t$. Wedge multiplication and the cone map
$-\int_0^1$ give the asserted product. This is an absolute product
regulator calculation, requiring no regulator on a relative lift.

Consider the actual class
\begin{equation}\label{eq:twodimkappa}
 \kappa=\sum_{g\in\F_5^\times}
             (\sigma_g\beta|_E)\,\eta_g\in K_3(\mathcal Q).
\end{equation}
At the fixed embedding write $b_g$ for the unique purely imaginary
representative of the point regulator of $\sigma_g\beta$.
These satisfy $b_{-g}=-b_g$, and at least one is nonzero by Borel
injectivity applied to all embeddings. Thus $(b_g)$ is a nonzero
anti-invariant vector, and in particular is not diagonal.

The weight-three regulator of $\kappa|_{\mathcal Q_s}$ is
\[
 \left[2\pi i\sum_g b_g e_g|_{\mathcal Q_s}\right].
\]
The vector inside brackets is a nonzero real Betti class: the only
root relation is diagonal, and
Lemma~\ref{lem:twodimbettipersistence} preserves it. A nonzero real
class cannot lie in the imaginary subspace $H^2(U,\R(3))$.
Therefore $\kappa|_{\mathcal Q_s}$ has nonzero regulator, and has
infinite order, for every $s\in S$.

To place this use exactly in the regulator's arithmetic site, all
classes and any one chosen $s$ spread to a finite localization
$\OO_E[1/N]$, enlarging the integer $N$ to include 5. The four roots
and their differences are units there. The surface
$xv+P(y)=0$, with $y$ inverted, is smooth over that ring; its
principal open defined by $s$ is a regular separated scheme of finite
type over $\Z$. The construction in \cite{BunkeTamme} therefore
applies at that finite stage. Naturality as $N$ and $s$ increase
gives exactly the regulators used above. Evaluation at the chosen
complex component is a projection, not an average over embeddings.

Finally, invert $u$ and $y$ after the flat base change from $A$ to
$B$. The divisor $x=0$ becomes the disjoint union of the four roots
$y=a_g$. The coefficient homomorphism on the $g$th component is
$W\to E$, $\rho\mapsto\rho^g$. Flat base change and the projection
formula therefore give the exact equality
\[
 c|_{\mathcal T}
   =\sum_g (\sigma_g\beta|_E)
                       [\OO_{D_g\cap\mathcal T}]
   =\kappa|_{\mathcal T}.
\]
The regularity of $\mathcal T$ legitimizes the identification of
$G$-theory and $K$-theory in this calculation; no such identification
on the singular ring $B$ is asserted.

Filtered-localization continuity gives
\[
 K_3(\mathcal T)=\mathop{\mathrm{colim}}_{s\in S}K_3(\mathcal Q_s)
 \qquad\cite[Chapter~IV, Section~6.4]{Weibel}.
\]
If a nonzero integer killed $c$, it would kill the image of $\kappa$
in this colimit and hence kill its restriction at some finite stage.
That contradicts the nonzero regulator at every such stage. Thus $c$
has infinite order. Its generic image was zero by
\eqref{eq:twodimsupported}, proving
Theorem~\ref{thm:twodimgeometric}.

\begin{remark}[Scope]
The coefficient residue field is $\F_5(T)$, not a finite field, and
the source ring is uncompleted. No consequence for its completion
follows from this complex argument. The source is ramified:
$5\in\mathfrak m_A^4$. In fact it cannot be geometrically regular
over any DVR. Indeed, $5$ is prime in $A$: the special-fiber equation
$x^4+y^4+Txy^3$ is irreducible over $\F_5(T)$, since
$s^4+Ts+1$ is primitive and linear in $T$. For a local DVR map
with uniformizer $\pi$, the equality $5=\varepsilon\pi^e$ would
therefore force $e=1$ and $\pi A=5A$; this fiber is singular.
Any putative geometrically regular DVR presentation is local after
localizing the base at the image of the closed point.

There is also no coefficient-ring map $W\to A$, since
$\Frac A=\Q(x,y)$ has no nonrational algebraic constants. Thus a
coefficient-ring splitting cannot supply the lift required for the
elementary projection-formula vanishing $i_*i^*\alpha=0$.
\end{remark}

\begin{remark}[An explicit number-field coefficient]
The coefficient can be fixed using de Jeu's rational root-of-unity
class $[-\rho]_2\in K_3(E)_{\Q}$. Its wedge boundary is zero because
$-\rho$ is torsion. The all-embedding regulator formula
\cite[Theorem~2.3]{DeJeuCurves} identifies its complex regulator,
up to a fixed nonzero normalization, with $iD(-\rho)$, where $D$ is
the Bloch--Wigner dilogarithm. At $\rho=e^{2\pi i/5}$,
$D(-\rho)=-\operatorname{Cl}_2(3\pi/5)\ne0$.
Indeed, $\operatorname{Cl}_2(\theta)=
-\int_0^\theta\log(2\sin(t/2))\,dt$ is decreasing on
$(\pi/3,\pi)$ and is zero at $\pi$.
Clear one rational denominator, lift to $K_3(W)$, and apply
$1-\sigma_{-1}$. This gives an actual coefficient satisfying all
conditions of the proof. In particular the choice of coefficient
need not depend on a regulator existence argument alone.
\end{remark}

\section{The same ring in every higher odd degree}

\begin{theorem}[Proposed higher odd consequence]
\label{thm:twodimalloddgeometric}
For this single uncompleted ring $A$ and every $n\ge2$, the kernel of
$K_{2n-1}(A)\to K_{2n-1}(\Frac A)$ contains an actual infinite-order
class. Its rational dimension is at least two for even $n$ and at
least one for odd $n$.
\end{theorem}
\begin{proof}
Borel's ranks are $\dim_{\Q}K_{2n-1}(E)_{\Q}=2$ and
$\dim_{\Q}K_{2n-1}(\Q)_{\Q}=0$ for even $n$, respectively $1$
for odd $n$. Restriction and transfer satisfy
$\operatorname{res}\operatorname{Tr}=\sum_g\sigma_g$ and
$\operatorname{Tr}\operatorname{res}=4$. Thus the invariant subspace
is exactly the restriction from $\Q$, and its trace-zero complement
$V_n$ has the asserted dimension. Quillen localization gives
$K_{2n-1}(W)_{\Q}\simeq K_{2n-1}(E)_{\Q}$.
Choose an actual $\gamma\in K_{2n-1}(W)$ with noninvariant Borel
image and put $\beta_n=4\gamma-\sum_g\sigma_g\gamma$.
Its regulator is nonzero and its Galois sum is zero exactly.

Set $c_n=i_*\beta_n$. Flat base change again identifies its image
on $\mathcal T$ with $\sum_g(\sigma_g\beta_n|_E)\eta_g$.
At a complex embedding, represent the point regulators
$b_g\in\C/\R(n)$ in the complementary real line $i\R(n)$.
They form a nonzero vector with sum zero, hence a nonconstant vector.
For each finite open $U=\mathcal Q_s$, the product regulator is
represented by $2\pi i\sum_gb_ge_g$ in
\[
 \HD^3(U,\R(n+1))=H^2(U,\C)/H^2(U,\R(n+1)),
 \qquad n+1\ge3.
\]
Multiplication by $2\pi i$ induces an isomorphism
$\C/\R(n)\simeq\C/\R(n+1)$. The root lattice and
Lemma~\ref{lem:twodimbettipersistence} therefore make this class
nonzero. The preceding continuity argument gives infinite order
and generic vanishing. Applied to every nonzero vector of $V_n$,
the same argument injects $V_n$ into the rational kernel and proves
the bounds.
\end{proof}

The restriction $n\ge2$ is essential here: at $n=1$, the target
twist is two and $F^2$ on a complex affine surface is not zero.
The displayed quotient formula cannot be used in that degree.

% Optional independent supplement for the uncompleted geometric short proof.
% Uses only its Picard/Betti persistence and the complex BT regulator.
% No BDJ, syntomic, crystalline, or completed K-theory comparison is used.

\section{An explicit triple, detected by a complex period}
\label{sec:twodimexplicitgeometric}

\begin{theorem}[Proposed explicit uncompleted class]
\label{thm:twodimexplicitgeometric}
For the ring $A$ of Theorem~\ref{thm:twodimgeometric}, the actual
unit product
\begin{equation}\label{eq:twodimexplicitgeometric}
 c_{\rm exp}=\left\{2+y,\ 1-x,\
 1+(1-x)(x^3+Ty^3)y\bigl((1+y)^5+1\bigr)\right\}\in K_3(A)
\end{equation}
has infinite order and restricts to zero in $K_3(A[1/x])$ integrally.
\end{theorem}

\begin{proof}
Put $t=-1-y$, $v=x^3+Ty^3$, and
\[
 H(t)=(t^5-1)(t+1)=y\bigl((1+y)^5+1\bigr),\quad
 b=vH(t),\quad \lambda=1-x,\quad \mu=1+\lambda b.
\]
Since $xv=-\Phi_{10}(t)$, we have $1-xb=t^{10}$.
The entries $g=1-t=2+y$, $\lambda$, and $\mu$ have respective
residues $2,1,1$, and $t$ is a unit. With the convention
$\langle r,s\rangle=\{r,1-rs\}$ when $r$ is a unit,
the Dennis--Stein relations give
\[
 \delta=\langle x,b\rangle
 =\langle x,b+1-xb\rangle=\langle x,\mu\rangle
 =\{\lambda t^{10},\mu\}.
\]
Here $\langle x,1\rangle=0$ and $1-x\mu=\lambda t^{10}$;
see \cite[Chapter~III, Definition~5.11 and Exercise~5.6]{Weibel}.
Steinberg gives $\{1-t,t^{10}\}=0$, so
\begin{equation}\label{eq:twodimexplicitDS}
 [1-t]\delta=c_{\rm exp},\qquad
 [1-t]\delta|_{D(x)}=\{1-t,x,t^{10}\}=0.
\end{equation}

For nonvanishing consider the smooth affine surface
\[
 Q^0=\Spec E[X,V,t]/(XV+\Phi_{10}(t)),\qquad
 U=Q^0[1/(t(1-t))].
\]
On $U$, the Dennis--Stein product
$\alpha=[1-t]\langle X,VH(t)\rangle$ is defined and vanishes
on $U[1/X]$. The whole divisor $D=V(X)$ lies in $U$.
Localization and devissage give an actual supported lift
$\kappa^0\in K_3(Q^0)$ with $\kappa^0|_U=\alpha$.
There is no requirement that the last two entries of the triple
be units on the analytic cycle below; we use its equal
Dennis--Stein expression there.

\emph{Independence of the supported lift.}
This lift is unique rationally as a class on $Q^0$.
Here is a direct $K$-theoretic check, independent of its periods,
using localization and the Laurent fundamental theorem \cite{Weibel}.
Put $Z=\mathbf A^1_{E}\setminus\{0,1\}$.  Then
$U[1/X]\simeq\mathbf G_m\times Z$, and localization on the
affine line, split by the units $t$ and $1-t$, gives
\[
 K_r(Z)\simeq K_r(E)\oplus
       K_{r-1}(E)[t]\oplus K_{r-1}(E)[1-t].
\]
Together with the Laurent fundamental theorem this gives
\[
\begin{split}
 K_4(U[1/X])\simeq{}&
 K_4(E)\oplus K_3(E)[t]\oplus K_3(E)[1-t]\\
 &{}\oplus K_3(E)[X]
 \oplus K_2(E)[X][t]\oplus K_2(E)[X][1-t].
\end{split}
\]
The first three summands extend to $U$, so their boundary in
$K_3(D)=K_3(E)^4$ is zero.  For a number field,
$K_2(E)_{\mathbf Q}=K_4(E)_{\mathbf Q}=0$: localization
reduces this to the rational even $K$-group vanishing for its
number ring and the torsion odd $K$-groups of finite fields.
The only remaining boundary is that of
$\beta[X]$, namely the diagonal tuple
$\pm(\beta,\beta,\beta,\beta)$ with a common convention-dependent
sign, since $X$ has order one on each root component.  Hence
\[
 \ker\bigl(K_3(D)_{\mathbf Q}\longrightarrow K_3(U)_{\mathbf Q}\bigr)
       =\Delta K_3(E)_{\mathbf Q}.
\]
The diagonal pushes forward to zero on $Q^0$, by the
length-one free resolution of the principal divisor $X=0$.
Consequently any two actual supported lifts of $\alpha$ give
the same class in $K_3(Q^0)_{\mathbf Q}$.
Moreover, since
$Q^0[1/X]\simeq\mathbf G_m\times\mathbf A^1_{E}$, the same
Laurent calculation gives
$\ker(K_3(D)\to K_3(Q^0))=\Delta K_3(E)$ already integrally.
Thus its supported coefficients are unique modulo the diagonal.

Fix an embedding of $E$ into $\C$, put $\theta_0=\pi/5$,
$r=e^{i\theta_0}$, and take the chord
$t=\gamma(s)=\cos\theta_0+i\sin\theta_0\,s$, $-1\le s\le1$.
It joins the two primitive tenth roots $\bar r,r$ and avoids
$0,1,-1$ and all other roots. Its matching sphere is
\[
 X=\sqrt{|\Phi_{10}(t)|}\,e^{i\vartheta},\qquad
 V=-\frac{\Phi_{10}(t)}{|\Phi_{10}(t)|}\,\overline X,
\]
with $X=V=0$ at both ends, oriented by $ds\wedge d\vartheta$.
The displayed phase ensures $XV=-\Phi_{10}(t)$. This is smooth
at each pole since $|X|^2$ is a smooth positive multiple of its
distance in $s$ from the corresponding endpoint.

On one neighborhood of the chord choose the principal branches
\[
 L=\operatorname{Log}(1-t),\qquad F=\operatorname{Li}_2(t),
 \qquad dF=-L\,dt/t.
\]
They agree with the usual series inside the unit disc. The
weight-two holomorphic curvature of the Dennis--Stein class is
\[
 \omega=10\frac{dX}{X}\wedge\frac{dt}{t}
       =-\frac{dX\wedge d(VH)}{1-XVH},
\]
so it extends across both poles. The absolute multiplicative
interval model of \cite[Proposition~2.8, Theorem~2.31 and
Lemma~4.5]{BunkeTamme} represents $r_3(\alpha)$ on this
neighborhood by $[-L\omega]$. On the punctured sphere,
$L\omega=10d(F\,dX/X)$. Stokes' theorem, with the two smooth
caps shrunk to their poles, therefore gives
\begin{align}
 \int_{\Sigma}r_3(\alpha)
 &=-20\pi i\bigl(F(r)-F(\bar r)\bigr)\pmod{\R(3)}
       \label{eq:twodimexplicitperiod}\\
 &=40\pi D_{\rm BW}(e^{i\pi/5})>0.\notag
\end{align}
Here $\R(3)=i\R$; the last quantity is a nonzero real period.
Conjugation and $|r|=1$ give $F(r)-F(\bar r)=2iD_{\rm BW}(r)$.
Positivity follows directly from
$D_{\rm BW}(e^{i\theta})=-\int_0^\theta
\log(2\sin(v/2))\,dv$ for $0<\theta<\pi/3$.
No logarithm of $X$, or unspecified phase of $\Phi_{10}(t)$,
enters this calculation.

The substitutions $X=x$, $V=x^3+Ty^3$, $t=-1-y$ identify
$Q^0[1/(t+1)]$ with $\mathcal Q$ in
\eqref{eq:twodimdanielewski}. The entire sphere lies on this open.
Write $\kappa=\kappa^0|_{\mathcal Q}$. Since every component of
$D$ is an affine line, homotopy invariance and the projection
formula express $\kappa=\sum_g\gamma_g\eta_g$ with
$\gamma_g\in K_3(E)$. If $b_g\in i\R$ represents its point
regulator, multiplicativity represents $r_3(\kappa)$ by the
\emph{real root class}
\[
 v_\kappa=2\pi i\sum_g b_g e_g\in H^2(\mathcal Q(\C),\R).
\]
The period \eqref{eq:twodimexplicitperiod} proves $v_\kappa\ne0$.
Lemma~\ref{lem:twodimbettipersistence} keeps this real class
nonzero on every $\mathcal Q_s$, so its image in
\eqref{eq:twodimdeligne} is nonzero at every such stage.
As above, each use of the regulator spreads to a regular finite
type model over a finite localization of the number ring.

Finally $t$ and $1-t$ are units on $\mathcal T$, and
\eqref{eq:twodimexplicitDS} gives the exact equality
$\kappa|_{\mathcal T}=c_{\rm exp}|_{\mathcal T}$.
If a nonzero integer killed $c_{\rm exp}$, filtered-localization
continuity would kill $\kappa|_{\mathcal Q_s}$ at some finite
stage, contrary to its regulator. This proves infinite order.
The argument asserts no consequence for completion.
\end{proof}

% New uncompleted generalization, proposed and under internal review.
% Uses the standard inputs already cited in the geometric short manuscript.
% This fragment does not assert any completed analogue.

\begin{corollary}[The same explicit symbol in ordinary Milnor theory]
\label{cor:twodimexplicitMilnor}
For the quartic local ring $A$ above, the ordinary Milnor symbol
\[
 m=\left\{2+y,\ 1-x,\
 1+(1-x)(x^3+Ty^3)y\bigl((1+y)^5+1\bigr)\right\}
       \in K_3^M(A)
\]
has infinite order and has zero image in
$K_3^M(\Frac A)$, integrally. In particular the ordinary
Milnor Gersten augmentation also fails after rationalization.
\end{corollary}

\begin{proof}
Write
\begin{align*}
 t&=-1-y,\qquad g=1-t,\qquad \lambda=1-x,\qquad q=t^{10},\\
 b&=(x^3+Ty^3)y\bigl((1+y)^5+1\bigr),\qquad \mu=1+\lambda b.
\end{align*}
The defining equation gives $1-xb=q$ and hence
$1-x\mu=\lambda q$. The elements $g,\lambda,\mu,q$
are units in $A$: their residues are respectively $2,1,1,1$.
In its fraction field, $x$ and $x\mu$ are also nonzero.
The two Steinberg relations and skew-symmetry in field
Milnor $K$-theory give
\begin{align*}
 0&=\{x\mu,\lambda q\},\qquad \{x,\lambda\}=0,\\
 \{x,q\}&=-\{\mu,\lambda\}-\{\mu,q\}
            =\{\lambda q,\mu\}.
\end{align*}
Furthermore $\{g,q\}=10\{1-t,t\}=0$. Therefore, without
rationalizing,
\[
 m=\{g,\lambda,\mu\}
   =\{g,\lambda q,\mu\}
   =\{g,x,q\}
   =-\{g,q,x\}=0
       \quad\text{in }K_3^M(\Frac A).
\]
The natural homomorphism $K_3^M(A)\to K_3(A)$ induced
by products of units sends $m$ to the already detected
infinite-order class $c_{\rm exp}$. Thus no nonzero integer
can kill $m$. This only uses the natural product homomorphism;
it does not use a comparison isomorphism or a retraction.
\end{proof}

\section{A tame family at every residue prime}
\label{sec:twodimtamefamily}

\begin{theorem}[Proposed family in local dimension two]
\label{thm:twodimtamefamily}
Let $p$ be any prime and let $d\ge3$ satisfy $p\nmid d$. Put
\[
 A_{p,d}=\left(\Z_{(p)}[T,x,y]/
       (p+x^d+y^d+Txy^{d-1})\right)_{(p,x,y)}.
\]
This is a regular local ring of dimension two, with residue field
$\F_p(T)$, and
\[
 \dim_{\Q}\ker\bigl(K_3(A_{p,d})_{\Q}
       \longrightarrow K_3(\Frac A_{p,d})_{\Q}\bigr)
       \ \ge\ \lfloor d/2\rfloor.
\]
In particular these dimensions are unbounded with the residue
prime and local dimension fixed. The theorem concerns the
uncompleted rings displayed above.
\end{theorem}

\begin{proof}
\emph{The actual source and the flat splitting ring.}
Write $A=A_{p,d}$. Its regular ambient local ring has parameters
$p,x,y$, and its defining equation has linear term $p$.
Thus $A$ is regular of dimension two with maximal ideal $(x,y)$.
Solving the equation for $T$ gives $\Frac A=\Q(x,y)$.

Choose $\theta^d=-p$ and set
\[
 K=\Q(\theta),\qquad W_0=\Z_{(p)}[\theta],\qquad
 A/(x)=W_0[T]_{(\theta)}.
\]
Eisenstein's criterion gives $[K:\Q]=d$; $W_0$ is a DVR with
uniformizer $\theta$ and residue field $\F_p$.
Quillen localization gives
$K_3(W_0)_{\Q}\simeq K_3(K)_{\Q}$, and every actual
$K_3(K)$ class lifts to $K_3(W_0)$, since $K_2(\F_p)=0$.
Pull back such a coefficient $\beta$ to $A/(x)$ and push it
forward to obtain $c_\beta\in K_3(A)$. Its restriction to
$A[1/x]$ is zero. We will prove that the resulting rational map
$K_3(K)_{\Q}\to K_3(A)_{\Q}$ is injective.

Let $E=K(\mu_d)$ and choose a prime $\mathfrak P$ above $p$.
Put $W=(\mathcal O_E)_{\mathfrak P}$ and let $k$ be its finite
residue field. The completion of $E$ at this prime is
$\Q_p(\mu_d,\theta)$. Since $p\nmid d$, adjoining $\mu_d$
is unramified, and $Z^d+p$ remains Eisenstein over that
unramified extension. Consequently $\theta$ is a uniformizer
of $W$, and the $d$ roots of unity $\zeta_i\in W$ have distinct
reductions in $k$.
The ring $W$ is flat over $\Z_{(p)}$, but need not be finite.
Use the explicit local ring
\[
 B=(A\otimes_{\Z_{(p)}}W)_{(\theta,x,y)}
  =\left(V[x,y]/
      (x^d+y^d+Txy^{d-1}-\theta^d)\right)_{(\theta,x,y)},
 \qquad V=W[T]_{(\theta)}.
\]
The ideal indicated is maximal with residue field $k(T)$, and
$A\to B$ is flat by tensor product followed by localization.
No finite-flat assertion about this map is used.

\emph{The exceptional curve and its root points.}
The projectivized tangent cone of $B$ is
\[
 C:\ X^d+Y^d+TXY^{d-1}=U^d
       \subset\mathbf P^2_{k(T)}.
\]
It is geometrically smooth. A singular point must have $U=0$;
then $Y\ne0$ and $q(s)=s^d+Ts+1$ would have a multiple root.
The identity $dq-sq'=(d-1)Ts+d$ shows this can occur only at
finitely many algebraic values of $T$, or never if $p\mid d-1$.
Thus it cannot occur over $k(T)$.

The smooth plane curve $C$ is geometrically integral, so the
associated graded ring of $B$ is a domain. Krull intersection
then shows that $B$ is a domain. The blow-up
$\mathcal Y=\operatorname{Bl}_{(\theta,x,y)}\Spec B$
is regular and has the single reduced exceptional curve $C$.
Here are the local details needed for this assertion. Its
$\theta$-chart has equation
$X^d+Y^d+TXY^{d-1}=1$, with $x=\theta X$, $y=\theta Y$;
the other two strict-transform charts reduce to the corresponding
charts of the smooth $C$. A nonzero differential tangent to
the ambient exceptional plane therefore proves regularity at
every point of $C$. The nonregular locus is closed by excellence.
Its proper image in the local scheme $\Spec B$, if nonempty,
would contain the closed point, a contradiction. This proves
regularity everywhere.
Off the closed point the blow-up is an isomorphism; the
hypersurface $S_2$ property then also shows $B$ is normal.

The tautological ideal gives
$\mathcal O_{\mathcal Y}(C)|_C=\mathcal O_C(-1)$.
The root primes $(x,y-\theta\zeta_i)$ have strict transforms
meeting $C$ at $P_i=[1:0:\bar\zeta_i]$, each with multiplicity
one. Indeed, on the $\theta$-chart the $Y$-derivative is a unit
at $X=0,Y=\zeta_i$, and the two regular parameters there are
$\theta,X$.

\emph{The exceptional Picard lattice.}
On $C^*=C\cap\{UY\ne0\}$ the root-point classes have only
the diagonal relation. To prove this, extend $k$ to $\bar k$ and
consider
\[
 M=\Spec\bar k[X,Z,Y^{\pm1}]/(XZ-(1-Y^d)),\qquad
 T=(Z-X^{d-1})/Y^{d-1}.
\]
On inverting $X$ its ring is the Laurent UFD
$\bar k[X^{\pm1},Y^{\pm1}]$. Divisor localization and its units
$cX^aY^b$ give
$\Pic(M)=\Z^d/\Z(1,\ldots,1)$, generated by
$(X,Y-\bar\zeta_i)$.

Every finite fiber $T=t$ is geometrically integral. Under
$s=X/Y$, $w=1/Y$ it is
$w^d=s^d+ts+1$, $w\ne0$.
The polynomial $q_t=s^d+ts+1$ has $\deg\gcd(q_t,q_t')\le1$
by $dq_t-sq_t'=(d-1)ts+d$ and $d\ne0$ in $\bar k$.
If it were an $\ell$th power for a prime $\ell\mid d$,
that degree would be at least $d-d/\ell\ge2$, a contradiction.
A rational-function power would already be a polynomial power.
Since $\bar k$ contains $\mu_d$ and $p\nmid d$, this also
proves irreducibility of $w^d-q_t$: the root extension is
cyclic, and a proper degree would express $q_t$ as such a power.
Each removed $T$-fiber is therefore one principal prime divisor,
of multiplicity one. Divisor localization gives
$\Pic(M)\simeq\Pic(M_{\bar k(T)})$.
Its generic fiber is $C^*_{\bar k(T)}$, proving the root-point
assertion.

\emph{Arithmetic Picard and Betti persistence.}
The generic affine surface and test localization are
\begin{align*}
 Q&=\Spec E[T,x,y,y^{-1}]/(p+x^d+y^d+Txy^{d-1})\\
  &\simeq\Spec E[x,v,y^{\pm1}]/(xv+y^d+p),\qquad
 \mathcal T=\Spec B[1/\theta,1/y],
\end{align*}
where $v=x^{d-1}+Ty^{d-1}$. Inverting $x$ again gives a
Laurent UFD. Thus $\Pic(Q)=\Z^d/\Z(1,\ldots,1)$, generated
by $D_i=(x,y-\theta\zeta_i)$; this remains true after every
coefficient field extension.
The map $\Pic(Q)\to\Pic(\mathcal T)$ is injective.
Indeed, a root combination made principal on $\mathcal T$
can acquire on $\mathcal Y$ only strict divisors over $\theta=0$
or $y=0$, and a multiple of $C$. The former miss $C^*$ and
$\mathcal O_C(-1)$ is trivial on $U\ne0$.
Restriction to $C^*$ gives the same combination of the $P_i$,
so all coefficients are equal, which was already the zero
relation on $Q$.

Write $\mathcal O(\mathcal T)=\varinjlim_s\mathcal O(Q_s)$,
where $s$ runs over the permitted arithmetic denominators.
Each $E$-prime divisor removed by such an $s$ has zero Picard
class by the injection just proved. Its geometric components
are one Galois orbit. The absolute Galois group over $E$ acts trivially on the torsion-free
root Picard lattice, so each geometric component also has zero
class. The degree-two support sequence consequently gives
\[
 H^2(Q(\C),\R)\lhook\joinrel\longrightarrow
 H^2(Q_s(\C),\R)\qquad\text{for every permitted }s.
\]
The root Chern classes $e_i$ span an injected copy of
$\R^d/\R(1,\ldots,1)$: in the Gysin sequence of the $D_i$,
the angular class of $x$ has residue the diagonal and that of
$y$ has residue zero. This proves precisely the needed real
root-class persistence.

\emph{Detection of every nonzero coefficient.}
Let $\iota_i:K\hookrightarrow E$ send $\theta$ to
$\theta\zeta_i$, and put $\eta_i=1-[\mathcal O_Q(-D_i)]$.
Flat base change for the actual Gysin class and the projection
formula, followed by restriction to the regular $\mathcal T$,
give
\[
 c_\beta|_{\mathcal T}
    =\left.\sum_i\iota_i(\beta)\eta_i\right|_{\mathcal T}.
\]
At one complex embedding of $E$, the point regulators
$b_i\in\C/\R(2)$ run through all embeddings of $K$.
Choose their representatives in $i\R$. If $\beta\ne0$ in
$K_3(K)_{\Q}$, Borel injectivity makes this vector nonzero.
Conjugate embeddings have opposite coefficients and real
embeddings have coefficient zero, so it is not diagonal.

The absolute multiplicative BT regulator of the supported sum
is the real root class $2\pi i\sum_i b_i e_i$, modulo
$H^2(Q_s(\C),\R(3))$. It is nonzero on every permitted finite
open by the preceding injection. Here the analytic Deligne
target is $H^2(Q_s,\C)/H^2(Q_s,\R(3))$ because its complex
dimension is two and $F^3=0$; a nonzero real class survives this
quotient. Each individual regulator calculation spreads to a
regular finite type arithmetic model after inverting one
integer. The standard continuity
$K_3(\mathcal T)=\varinjlim_s K_3(Q_s)$ now rules out any
nonzero rational coefficient in the kernel of the Gysin map.

Finally $Z^d+p$ has one real root when $d$ is odd and none
when $d$ is even. Borel's rank is therefore
$\dim_{\Q}K_3(K)_{\Q}=r_2(K)=\lfloor d/2\rfloor$.
Choosing actual lifts of a rational basis yields the asserted
number of independent infinite-order generic-kernel classes.
All uses of localization, base change and regulators are the
standard inputs cited above \cite{Quillen,QuillenFinite,Borel,BunkeTamme,Weibel}.
\end{proof}

For example, $(p,d)=(2,3)$ gives
$\bigl(\Z_{(2)}[T,x,y]/(2+x^3+y^3+Txy^2)\bigr)_{(2,x,y)}$.
For any fixed $p$, arbitrarily large odd $d$ prime to $p$ give
the unbounded rank assertion without enlarging the residue
field $\F_p(T)$. No assertion about these rings after completion
is made by this complex argument.

\begin{remark}[A negative control at degree two]
The condition $d\ge3$ is used in the vertical-fiber calculation.
When $d=2$ and $p\ne2$, the fiber at $t=2$ has
$s^2+2s+1=(s+1)^2$, so the Kummer irreducibility argument is
unavailable and the vertical fiber splits. Localization can then
kill a root divisor difference. The proof therefore supplies no
degree-two analogue of the displayed tame family.
\end{remark}


% Optional uncompleted corollary to thm:twodimtamefamily.
% No real scalar proportionality or completed coefficient assertion is used.

\begin{corollary}[Proposed higher odd bounds for the tame family]
\label{cor:twodimtameallodd}
For $A_{p,d}$ as in Theorem~\ref{thm:twodimtamefamily} and every
$n\ge2$, the rational kernel of
$K_{2n-1}(A_{p,d})\to K_{2n-1}(\Frac A_{p,d})$ has dimension at least
\[
 \begin{cases}
 \lfloor d/2\rfloor,&n\text{ even},\\
 \lceil d/2\rceil-1,&n\text{ odd}.
 \end{cases}
\]
These bounds are realized by actual classes of infinite order.
For odd $d$ the bound is $(d-1)/2$ in every odd $K$-degree at least three.
\end{corollary}

\begin{proof}
Retain $K=\Q(\theta)$, $W_0$, the splitting field $E$, and the
$d$ embeddings $\iota_i:K\hookrightarrow E$ from the theorem.
Use the specified rational coefficient space
\[
 V_n=\ker\left(\operatorname{Tr}_{K/\Q}:
       K_{2n-1}(K)_{\Q}\longrightarrow K_{2n-1}(\Q)_{\Q}\right).
\]
For an actual $\gamma\in K_{2n-1}(K)$ set
$\beta=d\gamma-\operatorname{res}_{K/\Q}
                 \operatorname{Tr}_{K/\Q}\gamma$.
The identity $\operatorname{Tr}\operatorname{res}=d$ shows
that $\operatorname{Tr}\beta=0$ integrally. These actual
classes span $V_n$ rationally, because the displayed operator
acts as multiplication by $d$ on $V_n$.
Each $\beta$ lifts to $K_{2n-1}(W_0)$, since
$K_{2n-2}(\F_p)=0$. Localization, flat base change and the
projection formula give an actual generic-zero class whose
test image is $\sum_i\iota_i(\beta)\eta_i$.

At a fixed complex embedding of $E$ let $b_i$ be the point
regulators, represented in the complementary real line
$i\R(n)$ to $\R(n)$ in $\C$. Base change for field transfer
gives $\sum_i\iota_i(\beta)=
\operatorname{res}_{E/\Q}\operatorname{Tr}_{K/\Q}\beta=0$;
therefore $\sum_i b_i=0$. If $\beta$ is nonzero rationally,
Borel injectivity makes the vector $(b_i)$ nonzero and hence
nondiagonal. Its line-bundle product is
$2\pi i\sum_i b_i e_i$. The preceding Betti persistence
keeps this root class nonzero on every permitted finite open.
The relevant analytic Deligne group is
\[
 H^3_{\mathcal D,\mathrm{an}}(Q_s,\R(n+1))
   =H^2(Q_s,\C)/H^2(Q_s,\R(n+1)),
\]
because $n+1\ge3$ and $Q_s(\C)$ is a surface.
Multiplication by $2\pi i$ identifies
$\C/\R(n)$ with $\C/\R(n+1)$, so the product regulator
is nonzero. The same continuity argument proves injectivity
of the supported map on $V_n$; applying it to every rational
relation gives independent actual representatives.

Rational transfer is surjective, since its composition with
restriction is multiplication by $d$. Borel's ranks give
\[
 \dim V_n=
 \begin{cases}
 r_2(K),&n\text{ even},\\
 r_1(K)+r_2(K)-1,&n\text{ odd}.
 \end{cases}
\]
Here $r_2(K)=\lfloor d/2\rfloor$ and $r_1(K)$ is one for
odd $d$ and zero for even $d$, proving the asserted bounds.
\end{proof}

% Separate insertion for the uncompleted geometric manuscript.
% Depends on thm:twodimtamefamily and its Picard/Betti argument.
% No completed, henselized-local-ring, or p-adic comparison conclusion.

\section{One ring with infinite rank in every higher odd degree}
\label{sec:twodimstrictcoefficientinfinite}

\begin{theorem}[Proposed infinite-rank consequence]
\label{thm:twodimstrictcoefficientinfinite}
Fix a prime $p$ and an integer $d\ge3$ prime to $p$. Let
$V=\Z_{(p)}^{\mathrm{sh}}$, with residue field $\overline{\F}_p$,
and put
\[
 R=\left(V[T,x,y]/(p+x^d+y^d+Txy^{d-1})\right)_{(p,x,y)}.
\]
Then $R$ is a Noetherian regular local ring of dimension two,
with residue field $\overline{\F}_p(T)$. For every $n\ge2$,
\[
 \dim_{\Q}\ker\!\left(K_{2n-1}(R)_{\Q}
       \longrightarrow K_{2n-1}(\Frac R)_{\Q}\right)=\infty.
\]
The coefficient DVR $V$ is strictly henselian; no henselianity
of the displayed $R$ is asserted. This theorem concerns this
single uncompleted ring.
\end{theorem}

\begin{proof}
\emph{The coefficient and splitting rings.}
The strict henselization $V$ is a DVR with uniformizer $p$,
and $F=\Frac V$ is algebraic over $\Q$; see the Stacks Project,
\href{https://stacks.math.columbia.edu/tag/0AP3}{Tag 0AP3}
and the proof of
\href{https://stacks.math.columbia.edu/tag/09ET}{Tag 09ET}.
It is excellent because it is a DVR of fraction characteristic
zero
(\href{https://stacks.math.columbia.edu/tag/07QW}{Tag 07QW(4)}).
The simple roots of $Z^N-1$ in $\overline{\F}_p$ lift by
Hensel's lemma whenever $p\nmid N$, so $\mu_N\subset V$.
The ambient local ring defining $R$ has parameters $p,x,y$,
and the defining equation has initial term $p$. Thus $R$ has
the stated regularity, dimension and residue field.

Choose $\theta^d=-p$ and set
\[
 W=V[\theta],\qquad E_\infty=F(\theta),\qquad
 G=\operatorname{Gal}(E_\infty/F)=\{\sigma_\zeta:\zeta\in\mu_d\},
 \quad \sigma_\zeta(\theta)=\zeta\theta .
\]
Eisenstein shows that $W$ is a finite free DVR extension of
$V$, with uniformizer $\theta$, residue field $\overline{\F}_p$,
and degree $d$. Its displayed Galois action preserves $W$.
Moreover
\[
 D=\Spec(R/(x))=\Spec W[T]_{(\theta)}
\]
is regular. Pullback to $D$ followed by its Gysin map
$i_*:K_{2n-1}(D)\to K_{2n-1}(R)$ therefore constructs
actual classes vanishing on $R[1/x]$.

\emph{The geometric detector over $E_\infty$.}
Use the flat splitting map to
\[
 B=(R\otimes_VW)_{(\theta,x,y)},\qquad
 \mathcal T=\Spec B[1/\theta,1/y],
\]
and the smooth surface
\[
 Q=\Spec E_\infty[x,v,y^{\pm1}]/(xv+y^d+p),
 \qquad v=x^{d-1}+Ty^{d-1}.
\]
The geometric proof of Theorem~\ref{thm:twodimtamefamily}
applies over these coefficient rings. Here is the scope check.
The blow-up still has the single smooth exceptional curve
\[
 C:\ X^d+Y^d+TXY^{d-1}=U^d
       \subset\mathbf P^2_{\overline{\F}_p(T)}.
\]
The same three charts prove its regularity along $C$, and
excellence and properness prove regularity everywhere. The
root strict transforms meet it transversely at
$[1:0:\bar\zeta]$. On $C\cap\{UY\ne0\}$ their Picard
classes have the lattice $\Z^d/\Z(1,\ldots,1)$:
the auxiliary surface
$XZ=1-Y^d$, $Y\ne0$, has that lattice, and every finite
$T=(Z-X^{d-1})/Y^{d-1}$ fiber is integral by the
$dq-sq'$ argument. These calculations already take place
over $\overline{\F}_p$ and do not require finite residue
constants.

Consequently, restriction of divisors to this exceptional open
gives
\[
 \Pic(Q)=\Z^d/\Z(1,\ldots,1)
       \lhook\joinrel\longrightarrow\Pic(\mathcal T),
\]
with generators $D_\zeta=(x,y-\zeta\theta)$.
The Picard calculation is unchanged by coefficient field
extension. Each $E_\infty$-prime divisor removed by a permitted
finite localization $Q_s$ has zero Picard class. Its geometric
components form one Galois orbit and all have the same class,
since the absolute Galois group over $E_\infty$ fixes its root
generators and the lattice is
torsion-free. Each component class is therefore zero.
The support sequence gives, for every embedding
$\jmath:E_\infty\hookrightarrow\C$,
\begin{equation}
 H^2(Q_\jmath(\C),\R)
      \lhook\joinrel\longrightarrow H^2((Q_s)_\jmath(\C),\R).
 \label{eq:strictcoefficientbetti}
\end{equation}
The Gysin sequence with complement $(\mathbf G_m)^2$ identifies
the root subspace with $\R^d/\R(1,\ldots,1)$, since the
angular classes of $x$ and $y$ have residues respectively the
diagonal vector and zero. This argument uses no regulator on
an infinite arithmetic base.

\emph{Finite actual coefficient spaces of unbounded rank.}
For $N>2$ divisible by $d$ and prime to $p$, choose compatible
roots of unity in $V$ and set
\[
 L_N=\Q(\mu_N)\subset F,\qquad K_N=L_N(\theta)\subset E_\infty .
\]
The chosen $p$-adic place of $L_N$ is unramified over $\Q_p$.
Eisenstein there gives $[K_N:L_N]=d$, and
$\operatorname{Gal}(K_N/L_N)$ is the restriction of $G$. Both fields
are totally imaginary. For fixed $n\ge2$, Borel's ranks and
restriction--transfer identities give
\begin{equation}
 V_{N,n}=
 \ker\!\left(\sum_{\sigma\in G}\sigma:
       K_{2n-1}(K_N)_{\Q}\to K_{2n-1}(K_N)_{\Q}\right),
 \qquad
 \dim_{\Q}V_{N,n}=\frac{(d-1)\varphi(N)}2.
 \label{eq:strictcoefficientrank}
\end{equation}
Indeed, the invariant subspace is exactly the restriction of
$K_{2n-1}(L_N)_{\Q}$, because
$\operatorname{res}\operatorname{Tr}=\sum_\sigma\sigma$ and
$\operatorname{Tr}\operatorname{res}=d$.
The two Borel ranks are $d\varphi(N)/2$ and $\varphi(N)/2$
in every degree under consideration.

Let $W_N=(\mathcal O_{K_N})_{\mathfrak P_N}$ be the DVR
at the prime induced by $W$. It maps to $W$, and $G$ preserves
it: above the chosen $L_N$ place the extension is totally
ramified and has a unique prime. Finite-field localization
gives
\[
 K_{2n-1}(W_N)_{\Q}\simeq K_{2n-1}(K_N)_{\Q},
\]
as well as actual surjectivity to the field, since
$K_{2n-2}$ of the finite residue field is zero.
Thus $V_{N,n}$ defines a rational linear map into
$K_{2n-1}(R)_{\Q}$ by coefficient pullback and $i_*$.
Its elements have actual representatives after clearing
denominators. If exact integral trace zero is wanted, replace
an actual lift $\gamma$ by
$d\gamma-\sum_{\sigma\in G}\sigma\gamma$; this operation
acts as multiplication by $d$ on $V_{N,n}$.

\emph{Detection, with both continuity passages explicit.}
Suppose $0\ne\beta\in V_{N,n}$ mapped to zero in
$K_{2n-1}(R)_{\Q}$. Borel injectivity supplies an embedding
$j_0:K_N\hookrightarrow\C$ at which its point regulator
is nonzero. Extend $j_0$ to an embedding
$\jmath:E_\infty\hookrightarrow\C$, using that
$E_\infty/K_N$ is algebraic.
For $\sigma\in G$, let $b_\sigma$ denote the point regulator
at $j_0\sigma$, represented in the complementary real line
$i\R(n)$ to $\R(n)$ in $\C$. Then
\[
 b_1\ne0,\qquad \sum_{\sigma\in G}b_\sigma=0,
\]
so this vector is not diagonal. The detecting embedding is
chosen for this coefficient; no single embedding is required
to detect every element of $V_{N,n}$.

Put $\eta_\zeta=1-[\mathcal O_Q(-D_\zeta)]$. Flat Gysin
base change and the projection formula identify the image of
the supported class on $\mathcal T$ with the restriction of
\[
 \alpha=\sum_{\zeta\in\mu_d}(\sigma_\zeta\beta)\eta_\zeta
       \in K_{2n-1}(Q)_{\Q}.
\]
The assumed vanishing therefore makes $\alpha$ zero on
$\mathcal T$. First, $K$-theory continuity for the filtered
localization
$\mathcal O(\mathcal T)=\varinjlim_s\mathcal O(Q_s)$
gives a finite permitted denominator set $s$ for which
$\alpha|_{Q_s}=0$ rationally.
Second, $E_\infty$ is a filtered union of number fields.
All coefficients of these finitely many denominators, the
root ideals, and the chosen $K$-classes descend to a number
field $E_0\subset E_\infty$ containing $K_N$.
Continuity for this filtered coefficient extension implies
that the asserted vanishing already holds on the corresponding
finite open over some larger number field
$E_1\subset E_\infty$.

Apply the published absolute multiplicative regulator to that
finite stage, after spreading it to a regular finite-type
arithmetic model by inverting one integer. Evaluate at
$\jmath|_{E_1}$. Its complex open is exactly
$(Q_s)_\jmath$, and naturality gives the product class
\[
 2\pi i\sum_{\zeta\in\mu_d}b_{\sigma_\zeta}e_\zeta
 \quad\text{in}\quad
 H^3_{\mathcal D,\mathrm{an}}((Q_s)_\jmath,\R(n+1))
 =\frac{H^2((Q_s)_\jmath,\C)}
        {H^2((Q_s)_\jmath,\R(n+1))}.
\]
The equality holds because $n+1\ge3$ and the complex
dimension is two. Multiplication by $2\pi i$ identifies
$\C/\R(n)$ with $\C/\R(n+1)$. The non-diagonal root
vector and \eqref{eq:strictcoefficientbetti} make this
regulator nonzero, contradicting the finite-stage vanishing.
This proves injectivity on every $V_{N,n}$.
No commutation of analytic cohomology with an infinite
coefficient colimit has been used.

Finally, take for example $N=d^a$ as $a\to\infty$.
The dimensions in \eqref{eq:strictcoefficientrank} are
unbounded, while $R$ remains fixed. This proves infinite
rational rank in every degree $2n-1\ge3$.
All coefficient lifts, transfers and products used above are
actual Quillen $K$-theory operations
\cite{Quillen,QuillenFinite,Weibel,Borel,BunkeTamme}.
\end{proof}

For each $p$ one may fix $d=3$ if $p\ne3$, and $d=4$
if $p=3$. This gives one proposed ring at each residue prime
with the stated infinite-rank conclusion in all higher odd
degrees. The argument proves neither persistence under
completion nor persistence under henselization of this $R$.

\clearpage
{\small
\begin{thebibliography}{9}
\bibitem{Borel}
A.~Borel, \emph{Stable real cohomology of arithmetic groups},
Ann. Sci. \'Ecole Norm. Sup. (4) \textbf{7} (1974), 235--272;
Proposition~12.2.
\href{https://doi.org/10.24033/asens.1269}{doi:10.24033/asens.1269}.

\bibitem{BunkeTamme}
U.~Bunke and G.~Tamme, \emph{Multiplicative differential algebraic
\(K\)-theory and applications}, Annals of \(K\)-Theory
\textbf{1} (2016), no.~3, 227--258;
Remark~2.6, Theorem~2.31 and Section~6.
\href{https://doi.org/10.2140/akt.2016.1.227}{doi:10.2140/akt.2016.1.227}.

\bibitem{DeJeuCurves}
R.~de Jeu, \emph{Towards regulator formulae for curves over number
fields}, Compositio Math. \textbf{124} (2000), 137--194;
Theorem~2.3.
\href{https://arxiv.org/abs/math/9811192}{Author's version}.

\bibitem{Quillen}
D.~Quillen, \emph{Higher algebraic \(K\)-theory I}, in
\emph{Algebraic \(K\)-Theory I: Higher \(K\)-Theories},
Lecture Notes in Math. \textbf{341}, Springer, 1973, 85--147.
\href{https://doi.org/10.1007/BFb0067053}{doi:10.1007/BFb0067053}.

\bibitem{QuillenFinite}
D.~Quillen, \emph{On the cohomology and \(K\)-theory of the general
linear groups over a finite field}, Ann. of Math. (2)
\textbf{96} (1972), 552--586.
\href{https://doi.org/10.2307/1970825}{doi:10.2307/1970825}.

\bibitem{StacksBlowup}
The Stacks Project Authors, \emph{The Stacks Project},
Section~31.33, \emph{Blowing up}, standard charts and tautological
ideal identity.
\href{https://stacks.math.columbia.edu/tag/01OF}{Tag 01OF}.

\bibitem{Weibel}
C.~A.~Weibel, \emph{The \(K\)-book: An Introduction to Algebraic
\(K\)-theory}, Graduate Studies in Math. \textbf{145}, AMS, 2013;
Chapter~IV, Section~6.4.
\href{https://sites.math.rutgers.edu/~weibel/Kbook/Kbook.IV.pdf}{Author's Chapter IV}.
\bibitem{PriorIntegralK3}
OpenAI, \emph{An integral degree-three Gersten counterexample},
research preprint, 26 September 2026, p.~1.
\href{https://github.com/openai/math/blob/main/preprints/An-Integral-Degree-Three-Gersten-Counterexample-September-26-2026/paper.pdf}{Repository PDF}.

\bibitem{PriorIntegralK5}
OpenAI, \emph{An integral counterexample to Gersten's conjecture},
research preprint, 25 September 2026, p.~2.
\href{https://github.com/openai/math/blob/main/preprints/An-Integral-Counterexample-to-Gerstens-Conjecture-September-25-2026/An-Integral-Counterexample-to-Gerstens-Conjecture-September-25-2026.pdf}{Repository PDF}.

\end{thebibliography}}
\end{document}
